\begin{table}[!htbp]
\centering
\scriptsize
\setlength{\tabcolsep}{3pt}
\begin{tabular}{@{}llrrr@{}}
\toprule
Model & Comparison summary & Fit RMSE & Forecast MAE & Forecast check loss \\
\midrule
Q--DESN \(\AL\)--\(\RHS\)
  & Reported criterion-specific summary & 2.176 & 2.361 & \textbf{3.297} \\
  & Five-chain common specification & \textbf{1.993} & \textbf{2.323} & 3.309 \\
\addlinespace
Q--DESN \(\exAL\)--\(\RHS\)
  & Reported criterion-specific summary & 1.710 & 2.709 & 3.333 \\
  & Five-chain common specification & \textbf{1.396} & \textbf{2.473} & \textbf{3.319} \\
\bottomrule
\end{tabular}
\caption{Five-chain MCMC sensitivity analysis for the Gaussian simulation
family at \(p=0.25\) with 500 fitting observations. The criterion-specific rows
reproduce the corresponding summaries in the main article's Gaussian MCMC
table and are interpreted separately by criterion. Each five-chain row uses one
pre-specified model specification and reports criteria recomputed from
the coordinatewise median of five chain-specific posterior point paths. This
aggregation summarizes posterior point paths across independent chains.
Forecast criteria use the 1000-observation held-out period, with origins spaced
30 observations apart and horizons up to 30 steps ahead. Boldface
marks the lower value within each model and criterion. This sensitivity analysis
uses a point-path estimator distinct from the draw-wise posterior summaries.}
\label{qdesn:tab:supp-independent-mcmc-five-chain-sensitivity}
\end{table}
